\documentclass[11pt]{article}

\usepackage[margin=1in]{geometry}
\usepackage{amsmath,booktabs,array}
\usepackage[hidelinks]{hyperref}
\usepackage{xcolor}
\usepackage{enumitem}
\usepackage{amsfonts}

\setlist[itemize]{leftmargin=1.5em,itemsep=0.25em,topsep=0.35em}
\setlist[enumerate]{leftmargin=1.8em,itemsep=0.25em,topsep=0.35em}
\newcommand{\code}[1]{\texttt{#1}}
\newcommand{\status}[1]{\textcolor{blue!65!black}{\textbf{#1}}}

\title{Racer-LLM Experiment Progress Log}
\author{}
\date{Last updated: August 13, 2026}

\begin{document}
\maketitle

\section{Project direction}

The main goal is to study the behavior and limitations of LLM-based policies in
a real-time multi-agent racing task. The baseline work is meant to provide
stable and understandable comparisons,instead of an exact
reproduction of every low-level detail in earlier racing systems of $\alpha$-racer or IBR from Wang et. al.. Should not diverge and get stuck in off-topic details, like tuning for better LLM agent performance.

The current comparison plans will eventually include:
\begin{itemize}
  \item locally trained and released $\alpha$-RACER policies;
  \item an online iterated best-response (IBR) planner;
  \item a pure rule-based policy;
  \item a raw LLM policy;
  \item an LLM policy with explicit prompt rules.
\end{itemize}

The parameterized policies share the same local MPC and simulator wherever
possible. IBR is a separate online trajectory-optimization baseline and should
be described as such.

\section{August 10--12, 2026: local $\alpha$-RACER training}

\status{Status: implemented, trained, and provisionally validated.}

\subsection{$\alpha$-RACER in brief}

$\alpha$-RACER represents each car's high-level policy using five parameters
\[
  \theta_i=(q_i,\zeta_i,s_{1,i},s_{2,i},s_{3,i}).
\]
The parameters control how the MPC follows and modifies a precomputed
single-agent raceline, including tracking aggressiveness, speed along the
raceline, overtaking, and blocking behavior.

For a joint policy for three car $\theta=(\theta_0,\theta_1,\theta_2) \in \Theta^3$, car $i$ has a
long-term value
\[
  V_i(x,\theta)=\mathbb E\left[\sum_t\gamma^t r^i_t\right].
\]
In the paper they use artificial neuron network $\Phi$ to learn from generated racing records data and to approximate the changes of race state value respect to single agent variation:
\[
\begin{split}
  &\Phi(x,\theta'_i,\theta_{-i})-\Phi(x,\theta_i,\theta_{-i})\\
  &\qquad\approx
  V_i(x,\theta'_i,\theta_{-i})-V_i(x,\theta_i,\theta_{-i}).
\end{split}
\]
The function learned offline is later used in online policy searches for joint
parameters that give a high potential value and executes its own five-parameter
block.


\subsection{Local implementation}

The current implemented pipeline contains three
clearly separated stages.

\paragraph{1. Rollout collection.}
For each group, a baseline joint parameter set is sampled within the configured bounds of $(q,\zeta,s_1,s_2,s_3)$. Four races are simulated, resample the five control parameters of single agent:
\begin{enumerate}
  \item the baseline parameters for all three cars: $(\theta_0,\theta_1,\theta_2)$;
  \item a unilateral parameter change for car 0: $(\theta_0',\theta_1,\theta_2)$;
  \item a unilateral parameter change for car 1: $(\theta_0,\theta_1',\theta_2)$;
  \item a unilateral parameter change for car 2: $(\theta_0,\theta_1,\theta_2')$.
\end{enumerate}
Only one car's parameters change in each deviation race. This gives paired data
for learning the effect of unilateral policy changes.

At each control update, the code stores a 39-dimensional feature vector (relative progress, vehicle states, track curvature, and all 15 policy
parameters). The local single-step reward is currently the closing the gap or extending the lead against the fastest rivals, plus the safety term of off-road or collision penalty:
\[
  r^i_t=w_p\Delta s^i_t
  -w_o\mathbf{1}_{\mathrm{offroad},i}
  -w_c\mathbf{1}_{\mathrm{collision}}.
\]

The final feature data would be at the shape of ($N_{group}$, 4, time\_steps, 39). 4 represents 4 pairs of data per group and 39 represents the feature dimension as well as the size of the input for value function and potential function.

\paragraph{2. Value-function training.}
After a rollout finishes, discounted returns are computed backwards:
\[
  G^i_t=\sum_{k=t}^{T-1}\gamma^{k-t}r^i_k.
\]
And the value function $\widehat{V}_i(x,\theta)$ is trained using this Monte-Carlo learning to maximize the discounted accumulative returns $G_t$. The loss function:
\[
\mathcal{L}_{V_i}=\mathbf{MSE}(\widehat{V}_i(z_i),G^i_t)
\]

\paragraph{3. Potential-function training.}
The three trained value networks produce targets for unilateral value changes.
The potential network is trained so that
\[
  \widehat\Phi(x,\theta')-\widehat\Phi(x,\theta)
  \approx
  \widehat V_i(x,\theta')-\widehat V_i(x,\theta).
\]
Both sides are evaluated at the same baseline state. This avoids mixing a policy
change with the different state trajectory that the policy later causes.
The loss function:
\[
\mathcal{L}_{\Phi}=\mathbf{MSE}(\Delta \widehat V_i(x)-\Delta \widehat \Phi_i(x))
=
\frac{1}{\mathcal D_\Phi}\sum_{(x,\theta,\theta')}(\widehat V_i(x,\theta')-\widehat V_i(x,\theta),
\widehat\Phi(x,\theta')-\widehat\Phi(x,\theta))^2,
\]
Existing generated rollout data and trained checkpoints are reused automatically. If there is new rollout data the code would automatically re-train the later stage.

\subsection{Connection to the released code}

\begin{center}
\small
\begin{tabular}{p{0.20\linewidth}p{0.35\linewidth}p{0.35\linewidth}}
\toprule
Task & Released repository & Local project \\
\midrule
Collect races &
\code{collect\_data.bash}, \code{multi\_car\_blocking.py},
\code{merge\_data.py} &
\code{collect\_rollout\_data()} \\
Fit three values &
\code{train\_q\_model\_multi.py},
\code{train\_q\_model\_multi\_rel.py} &
\code{train\_value\_functions()} \\
Fit potential &
\code{train\_p\_model\_multi\_dist.py} &
\code{train\_potential\_function()} \\
Network definition &
\code{model\_arch.py::SimpleModel} &
\code{policy/alpha\_model.py::SimpleNeuralNetwork} \\
Online policy &
Released numerical simulator &
\code{policy/alpha\_racer.py::AlphaRacerPolicy} \\
\bottomrule
\end{tabular}
\end{center}

The local implementation keeps the original overall idea but exclude the ROS simulation dependency and change the chaining order of data pairs to more independent deviation groups of 4 shown previously.

\subsection{Training run recorded}

\begin{itemize}
  \item Track: Berlin.
  \item Vehicle model: local kinematic model.
  \item MPC horizon: 10.
  \item Data: 100 groups $\times$ 4 rollouts $\times$ 120 control samples.
  \item Rollout duration: 12 s.
  \item Discount: $\gamma=0.98$.
  \item Seed: 7.
  \item Value validation MSE: 21.43, 20.95, and 22.63.
  \item Best potential difference MSE: 0.0761.
\end{itemize}

\subsection{Initial three-position checkpoint comparison}

As an initial check, the local checkpoint and the released checkpoint
were each evaluated using the standard three-position sweep on the Berlin track.
Each race contained three homogeneous $\alpha$-RACER cars using the same
checkpoint and lasted 50 seconds. Therefore, each checkpoint was evaluated in
three races, producing nine vehicle trajectories in total.

\begin{center}
\small
\begin{tabular}{lcc}
\toprule
Checkpoint & Mean progress gain per car & Mean speed \\
\midrule
Released myopic MPC & 180.80 m & 3.56 m/s \\
Locally trained & 198.64 m & 3.90 m/s \\
\bottomrule
\end{tabular} 
\end{center}

In this initial comparison, the locally trained checkpoint increased the mean
progress gain by 17.85 m and the mean speed by approximately 0.35 m/s. Neither
configuration produced a recorded collision. The released checkpoint had no
recorded off-road time, while the local checkpoint accumulated 0.08 seconds of off-road time across the nine vehicle trajectories.

These results show that the locally trained checkpoint runs successfully and
achieves higher progress under the local vehicle dynamics. It is therefore
suitable for use as the current $\alpha$-RACER baseline. However, this is only
an initial check: the three races mainly permute the starting order and do not cover independent random seeds, different tracks, or a broad range of initial conditions.

\section{August 10--12, 2026: ordinary local IBR}

\status{Status: implemented and tested with SE disabled.}

The other thing we worked on was replacing the problematic IBR inherited from
the released $\alpha$-RACER code. The new version is based on the main idea in
Wang et al.: optimize one car against the latest predicted trajectory of the
other car, then optimize the other car against the updated first-car
trajectory, and repeat this for a small number of rounds. Only the first
control input is executed, and the complete procedure is solved again at the
next control update.

In short, one best-response update is
\[
  \theta_i^{\ell+1}
  \in\arg\min_{\theta_i} f_i(\theta_i,\theta_{-i}^{\ell}).
\]
The current solver is \code{policy/ibr/trajectory\_ibr.py::\_WangIBRSolver}, selected by
\code{IBR\_IMPLEMENTATION=local\_wang}. It does not copy Wang's low-level
solver. Instead, it uses the state and controls already used in this project:
\[
  x_i=[s_i,e_i,\Delta\psi_i,v_i]^\top,
  \qquad
  u_i=[\text{throttle}_i,\text{steer}_i]^\top.
\]
This seemed more sensible than reproducing Wang's polynomial trajectory
parameterization while running the resulting controls in a different local
simulator.

Wang's ordinary IBR mainly maximizes terminal track progress and puts vehicle
motion, track containment, and collision avoidance in the constraints. We
started from this minimal objective instead of copying the collection of error
terms in the old local IBR. After the tests described below, the cost currently
used is
\[
\begin{aligned}
J_i={}&-w_p(s_{i,N}-s_{i,0})
  +w_e\sum_{k=0}^{N-1}e_{i,k}^2\\
&+w_u\sum_{k=0}^{N-1}\|u_{i,k}-u_{\mathrm{neutral}}\|_2^2
  +w_{\Delta u}\sum_{k=0}^{N-1}\|u_{i,k}-u_{i,k-1}\|_2^2\\
&+\rho_t\|\epsilon_t\|_2^2
  +\rho_c\|\epsilon_c\|_2^2.
\end{aligned}
\]
The actual weights are $w_p=10$, $w_e=0.4$, $w_u=0.05$, and
$w_{\Delta u}=0.10$. There is no heading-error or reference-speed loss in the
final version. In other words, $e^2$ is the only additional state-tracking
term; it does not mean that the complete objective contains only $e^2$.

The $e^2$ term is not part of Wang's original objective; it is a small local
regularizer that discourages the optimizer from continuously using the edge of
the track. The IBR does not follow the precomputed raceline or directly track
its speed profile. Raceline speed is used only to define the nominal control in
the control-effort term.

The constraints use the shared local kinematic model and vehicle parameters,
normalized throttle and steering bounds, maximum speed, steering-rate limit,
and the same friction-based steering/speed limit used by the local MPC. Track
containment is written in Frenet coordinates. Collision avoidance uses a
smooth superellipse because the collision-free condition for two rectangles is
a logical OR and is inconvenient to place directly in IPOPT. The simulator
still checks actual contact using oriented vehicle rectangles.

Track and collision constraints have nonnegative slack variables with large
penalties. These slacks are not part of Wang's ideal hard-constrained problem.
We kept them because hard-constraint tests in the current discrete simulator
often became infeasible. For example, after a small boundary overshoot, the
next optimization may be unable to return inside the boundary in one step
while also satisfying steering-rate and friction limits. A similar problem can
happen when one car solves against a frozen opponent trajectory in a tightly
compressed pack. Once the solver fails, fallback control can move the car even
farther away and cause a failure loop. Slack magnitude and activation are now
saved in the NPZ and JSON results, so this numerical mechanism can be reported
rather than hidden.

\subsection{Tests used to choose the current objective}

The objective variants were compared in complete three-car races with different
starting positions, interaction conditions, and tracks. Progress alone was
stable but tended to use the track boundary too aggressively. Restoring the
older lateral, heading, and speed losses reduced this behavior but also made the
cars less flexible during interaction. Keeping only $0.4e^2$ gave the best
overall behavior: no nominal boundary violations or solver failures were seen
in this test set, while progress was not reduced. We therefore kept $e^2$ as a
small local regularizer and left the heading and speed losses disabled.

Hard track and collision constraints were also tested, but they caused repeated
infeasibility in difficult closed-loop situations. The high-penalty slack
variables were therefore retained. Some brief contact can still occur when all
three cars form a compressed pack, since the current solver follows Wang's
two-player structure and responds to the nearest opponent. This is treated as a
known limitation rather than addressed with more ad hoc loss terms.

The older \code{historical\_strict} implementation is still available for
comparison. It is treated as a historical local port of the released
$\alpha$-RACER IBR rather than the main baseline.

\section{August 13, 2026: initial LLM policy comparison}

We compared the current LLM, an LLM given explicit decision rules in its
prompt, and a deterministic policy executing the same rules. All three select
from the same four $\alpha$-RACER modes and use the same low-level MPC. Each
policy was tested from all three starting positions, first against two local
$\alpha$-RACER cars and then against two IBR cars.

\begin{table}[ht]
  \centering
  \small
  \begin{tabular}{lcccc}
    \toprule
    & \multicolumn{2}{c}{vs. local $\alpha$-RACER}
    & \multicolumn{2}{c}{vs. IBR} \\
    \cmidrule(lr){2-3}\cmidrule(lr){4-5}
    Policy & Wins & Progress (m) & Wins & Progress (m) \\
    \midrule
    Raw LLM          & 2/3 & 193.2 & 0/3 & 193.1 \\
    Rule-prompt LLM  & 0/3 & 111.3 & 0/3 & 161.6 \\
    Rule-only policy & 0/3 & 134.8 & 0/3 & 171.8 \\
    Local $\alpha$-RACER self-play & 1/3 & 193.0 & -- & -- \\
    \bottomrule
  \end{tabular}
  \caption{Mean ego progress over the three starting positions.}
\end{table}

The raw LLM was close to local $\alpha$-RACER in progress, whereas adding the
explicit rules made both the LLM and deterministic policy too conservative.
They selected \emph{stable} frequently, although it is not a dedicated recovery
controller. The raw LLM was faster but also less consistent in safety, with
occasional contact or off-road recovery.

LLM inference did not meet the one-second decision interval: raw prompts took
about 3--4 s and rule prompts about 4--5 s. Because the simulator pauses for the
response, these runs evaluate policy decisions but not the physical effect of
inference delay.

\section{How IBR will be used in the comparison}

IBR is not a like-for-like high-level policy. The other methods select MPC
parameters and retain the raceline speed profile; IBR directly selects throttle
and steering, maximizes progress, and may use the full $5\,\mathrm{m/s}$ limit.

\begin{table}[ht]
  \centering
  \small
  \begin{tabular}{lccc}
    \toprule
    Candidate against local $\alpha$-RACER & Wins & Progress (m) & Speed (m/s) \\
    \midrule
    IBR, global speed limit    & 3/3 & 237.3 & 4.74 \\
    IBR, raceline-speed ablation & 0/3 & 183.0 & 3.61 \\
    Local $\alpha$-RACER self-play & 1/3 & 193.0 & 3.91 \\
    \bottomrule
  \end{tabular}
  \caption{The IBR result is strongly affected by its different speed objective.}
\end{table}

The unrestricted IBR result therefore represents the performance of a strong
full-trajectory online optimizer, not proof that its interaction reasoning is
better than the LLM. The main like-for-like comparison remains raw LLM,
rule-prompt LLM, rule-only, and $\alpha$-RACER, which share the five-parameter
action space and low-level MPC. The raceline-speed IBR is kept only as an
ablation.

\end{document}
